\documentclass[11pt]{article}
\usepackage{Style}

\begin{document}

\begin{center}
    \Large\textsc{El Teorema de Hurewicz} \\
    \normalsize Benjamín Mateluna
\end{center}

\noindent Esta presentación tiene como objetivo hablar sobre el teorema de Hurewicz, cuyo 
enunciado establece que el primer grupo de homología es isomorfo a la abelianización del grupo 
fundamental; esto es, al grupo resultante de cocientar este último por su conmutador, que 
denotaremos por
\begin{equation*}
    \widetilde{\pi}_{1}(X,x):=\pi_{1}(X,x)^{ab}=\frac{\pi_{1}(X,x)}{[\pi_{1}(X,x),\pi_{1}(X,x)]}
\end{equation*}
\noindent El resultado sirve como puente entre la teoría de homología y de homotopía, lo que 
permite traducir un problema homotópico en uno esencialmente algebraico al expresarlo en grupos de 
homología.

\vspace{2mm}
\noindent Para ello, se supondrá que el espacio topológico es arcoconexo, ya que este segundo 
invariante homotópico -- el grupo fundamental -- no es capaz de detectar las componentes 
arcoconexas, a diferencia de los grupos de homología. El morfismo a definir, digamos,
\begin{equation*}
    \phi:\pi_{1}(X,x)\to H_{1}(X)
\end{equation*}
es natural: dado un lazo basado en un punto, su imagen será el mismo lazo, pero visto ahora como 
una cadena en el complejo de cadenas de nuestro espacio.

\vspace{2mm}
\noindent Necesitaremos tres lemas para demostrar que dicha asignación está bien definida, que 
desciende al primer grupo de homología y que efectivamente es un morfismo de grupos. Una vez 
definido, dado que el primer grupo de homología es abeliano, este morfismo se factorizará a través 
de la abelianización del grupo fundamental debido a la propiedad universal del conmutador. Nuestra 
misión entonces será probar que tal homomorfismo, en realidad, es un isomorfismo.

\vspace{2mm}
\noindent Con este propósito, definiremos un morfismo auxiliar
\begin{equation*}
    \psi_{*}:H_{1}(X)\to\widetilde{\pi}_{1}(X,x)
\end{equation*}
\noindent y probaremos que corresponde a la inversa de nuestro candidato a isomorfismo. Se 
probarán dos lemas -- uno para cada composición a verificar -- y, para finalizar, se verán 
aplicaciones de este teorema, principalmente para calcular el primer grupo de homología a partir 
del grupo fundamental y viceversa.

\begin{thebibliography}{1}
    \bibitem{Bredon}
        \textsc{Bredon, Glen E.} (1993). 
        \textit{Topology and Geometry}, 1st ed. Springer.
\end{thebibliography}

\end{document}